\documentclass[11pt]{article}
\usepackage[a4paper,margin=22mm]{geometry}
\usepackage{fontspec}
\setmainfont{DejaVu Serif}
\setsansfont{DejaVu Sans}
\usepackage{microtype}
\usepackage{xcolor}
\usepackage{hyperref}
\usepackage{enumitem}
\usepackage{parskip}
\definecolor{Accent}{HTML}{971D32}
\definecolor{Muted}{HTML}{666666}
\hypersetup{colorlinks=true,urlcolor=Accent,linkcolor=Accent}
\setlist{leftmargin=1.6em,itemsep=0.3em,topsep=0.35em}
\setcounter{tocdepth}{2}
\renewcommand{\contentsname}{Contents}
\newcommand{\sectionrule}{\par\vspace{-0.5em}\noindent\textcolor{Accent}{\rule{\linewidth}{0.6pt}}\par}
\begin{document}

\begin{center}
{\sffamily\bfseries\color{Accent} TOEFL WORD OF THE DAY\par}
\vspace{8mm}
{\fontsize{42}{48}\selectfont\bfseries subjective\par}
\vspace{2mm}
{\Large\sffamily\color{Accent} /səbˈdʒɛktɪv/\par}
\vspace{2mm}
{\small\sffamily Local audio phrase: subjective\par}
{\small\sffamily Audio file: audios/subjective.mp3\par}
\vspace{2mm}
{\large\sffamily\color{Muted} adjective\par}
\vspace{5mm}
{\sffamily based on personal judgment \textbullet\ existing in personal experience\par}
\vspace{6mm}
{\large Subjective describes judgments or experiences shaped by an individual's feelings, opinions, or personal perception.\par}
\vspace{6mm}
\begin{minipage}{0.84\textwidth}
\itshape “Judgments about artistic quality are often highly subjective.”
\end{minipage}\par
\vspace{10mm}
\textbf{Date:} August 8th 2026\quad\textbullet\quad \textbf{Level:} B1--mid-B2 TOEFL
\vfill
Created and compiled by \textbf{Amir Kooshky}\par
\vspace{2mm}
\href{https://t.me/KooshkyTOEFL}{Telegram Channel}\quad\textbullet\quad
\href{https://www.instagram.com/kooshkytoefl}{Instagram}\quad\textbullet\quad
\href{https://t.me/KooshkyTOEFL\_pv}{Personal Telegram}
\end{center}
\newpage
\tableofcontents
\newpage

\section{Meanings and usage}
\sectionrule
\subsection{Meaning 1: Based on personal judgment}
\textbf{Part of speech:} adjective

\textbf{IPA:} /səbˈdʒɛktɪv/

\textbf{Pronunciation phrase:} subjective

\textbf{Local audio file:} audios/subjective.mp3

\textbf{Register:} neutral to formal; common in academic, professional, artistic, and everyday discussions

\textbf{Definition:} based on a person's opinions, feelings, or preferences rather than only on facts that can be measured or proved

\textbf{In simpler words:} based on personal opinion

\textbf{Typical pattern:} subjective + judgment, opinion, assessment, interpretation, or criterion

\subsubsection{Natural examples}
\begin{enumerate}
  \item Judgments about artistic quality are often highly subjective.
  \item The interview process was criticized because the final decisions seemed too subjective.
  \item Whether a building is beautiful is partly a subjective question.
  \item The researcher tried to separate subjective impressions from measurable evidence.
  \item Teachers should use clear criteria to make essay grading less subjective.
  \item Her review provides a subjective interpretation of the film rather than a neutral summary.
  \item Ideas about the ideal work environment are subjective and may vary considerably from person to person.
\end{enumerate}

\subsubsection{Nuance and implication}
\textbf{Central nuance:} The word emphasizes that a judgment is influenced by the person making it.

\textbf{Usual implication:} A subjective judgment is not automatically unreasonable or wrong, but different people may reach different conclusions.

\textbf{Compared with objective:} Subjective judgments depend on personal views or feelings. Objective judgments are based as much as possible on observable facts and agreed standards.

\subsubsection{Grammar notes}
\textbf{How to use it:} Use it before a noun, as in a subjective assessment, or after a linking verb, as in The assessment is subjective.

\textbf{What can follow it:} It commonly appears before nouns such as judgment, opinion, view, assessment, interpretation, impression, and criterion.

\textbf{Position in a sentence:} Words such as highly, largely, partly, inherently, and necessarily often come before subjective.

\subsubsection{Useful patterns}
\textbf{Pattern:} a subjective judgment or assessment

\textbf{Use:} Use this for a decision influenced by the evaluator's personal views.

\textbf{Example:} The selection process depends partly on a subjective assessment of leadership potential.

\textbf{Pattern:} be highly or largely subjective

\textbf{Use:} Use this when personal opinion strongly affects a judgment.

\textbf{Example:} The choice of the best design is highly subjective.

\textbf{Pattern:} a subjective interpretation of + noun

\textbf{Use:} Use this for an explanation shaped by one person's perspective.

\textbf{Example:} The article offers a subjective interpretation of the historical evidence.

\textbf{Pattern:} make something less subjective

\textbf{Use:} Use this when rules or measurements reduce the influence of personal opinion.

\textbf{Example:} A detailed scoring guide can make the evaluation less subjective.

\subsubsection{Common wrong usage}
\paragraph{Correction 1}
\textbf{Wrong:} The decision was subjective from the manager's personal opinion.

\textbf{Correct:} The decision was based on the manager's subjective opinion.

\textbf{Why:} Use subjective directly before the noun it describes, or say that something is based on a subjective judgment or opinion.

\paragraph{Correction 2}
\textbf{Wrong:} Subjective always means unfair or incorrect.

\textbf{Correct:} Subjective means influenced by personal opinion or feeling, but it does not always mean unfair or incorrect.

\textbf{Why:} A judgment can be subjective without being dishonest or unreasonable.

\paragraph{Correction 3}
\textbf{Wrong:} The results were very subjectively.

\textbf{Correct:} The results were very subjective.

\textbf{Why:} Use the adjective subjective after were. Subjectively is an adverb.

\subsection{Meaning 2: Existing in personal experience}
\textbf{Part of speech:} adjective

\textbf{IPA:} /səbˈdʒɛktɪv/

\textbf{Pronunciation phrase:} subjective

\textbf{Local audio file:} audios/subjective.mp3

\textbf{Register:} formal; common in psychology, medicine, sociology, philosophy, and research writing

\textbf{Definition:} experienced, perceived, or understood within an individual person's mind and therefore not fully observable or measurable by others

\textbf{In simpler words:} personally experienced or perceived

\textbf{Typical pattern:} subjective + experience, perception, feeling, report, rating, or well-being

\subsubsection{Natural examples}
\begin{enumerate}
  \item Pain is a subjective experience that can be difficult to measure precisely.
  \item The survey asked participants to report their subjective sense of well-being.
  \item People may give different descriptions of the same event because perception is subjective.
  \item The study combined biological measurements with patients' subjective reports of fatigue.
  \item Time can feel slower or faster depending on a person's subjective experience.
  \item The taste of the medicine was evaluated through subjective ratings provided by volunteers.
  \item Researchers distinguished objective sleep duration from subjective sleep quality.
\end{enumerate}

\subsubsection{Nuance and implication}
\textbf{Central nuance:} This sense focuses on what a person experiences internally, such as pain, satisfaction, taste, or perceived quality.

\textbf{Usual implication:} The experience may be real and important even when another person cannot directly observe or measure it.

\textbf{Compared with objective:} A subjective measure records a person's own experience or report. An objective measure records something observable or measurable independently, such as temperature, distance, or heart rate.

\subsubsection{Grammar notes}
\textbf{How to use it:} Use it before a noun naming an internal experience, perception, or self-reported measure.

\textbf{What can follow it:} Common nouns include experience, perception, feeling, report, rating, well-being, awareness, and quality.

\textbf{Position in a sentence:} In research writing, subjective is often contrasted directly with objective.

\subsubsection{Useful patterns}
\textbf{Pattern:} subjective experience of + noun

\textbf{Use:} Use this for how an event or condition feels to an individual.

\textbf{Example:} The interviews explored each patient's subjective experience of recovery.

\textbf{Pattern:} subjective report or rating

\textbf{Use:} Use this for information supplied by a person about their own experience.

\textbf{Example:} The researchers collected subjective ratings of stress after each task.

\textbf{Pattern:} subjective sense of + noun

\textbf{Use:} Use this for a personally felt state or impression.

\textbf{Example:} Regular social contact improved participants' subjective sense of belonging.

\textbf{Pattern:} subjective and objective measures

\textbf{Use:} Use this when comparing self-reported experience with independently observable data.

\textbf{Example:} The study used both subjective and objective measures of sleep quality.

\subsubsection{Common wrong usage}
\paragraph{Correction 1}
\textbf{Wrong:} Because pain is subjective, it is not real.

\textbf{Correct:} Pain is subjective because it is experienced personally, but it is still real.

\textbf{Why:} Subjective means internally experienced, not imaginary or unimportant.

\paragraph{Correction 2}
\textbf{Wrong:} The study measured subjectively well-being.

\textbf{Correct:} The study measured subjective well-being.

\textbf{Why:} Use the adjective subjective before the noun well-being.

\paragraph{Correction 3}
\textbf{Wrong:} The patient's subjective was recorded in the survey.

\textbf{Correct:} The patient's subjective experience was recorded in the survey.

\textbf{Why:} Subjective is an adjective and normally needs a noun after it in this position.

\section{Word family}
\sectionrule
\subsection{subjectivity}
\textbf{Part of speech:} uncountable noun

\textbf{IPA:} /ˌsʌbdʒɛkˈtɪvəti/

\textbf{Definition:} the influence of personal opinions, feelings, or experiences on judgment or perception

\textbf{Relationship to the headword:} This noun names the quality or condition of being subjective.

\textbf{Grammar pattern:} subjectivity in + judgment, assessment, interpretation, or experience

\textbf{Usage note:} It is normally used without a or an.

\subsubsection{Examples}
\begin{enumerate}
  \item The scoring system was designed to reduce subjectivity.
  \item The article discusses the role of subjectivity in historical interpretation.
\end{enumerate}

\subsection{subjectively}
\textbf{Part of speech:} adverb

\textbf{IPA:} /səbˈdʒɛktɪvli/

\textbf{Definition:} according to personal opinion, feeling, or experience rather than only observable facts

\textbf{Relationship to the headword:} This adverb describes a judgment or experience shaped by a person's own perspective.

\textbf{Grammar pattern:} judge, evaluate, experience, or interpret + subjectively

\subsubsection{Examples}
\begin{enumerate}
  \item Participants rated the difficulty of the task subjectively.
  \item The two paintings may be judged subjectively in different ways.
\end{enumerate}

\section{Common collocations}
\sectionrule
\subsection{highly subjective}
\textbf{Applies to:} Based on personal judgment

\textbf{Meaning:} strongly dependent on personal opinion or preference

\textbf{Pattern:} be highly subjective

\textbf{Example:} Deciding which novel is the most influential is highly subjective.

\subsection{subjective judgment}
\textbf{Applies to:} Based on personal judgment

\textbf{Meaning:} a decision influenced by the evaluator's personal views

\textbf{Example:} The final score still involves some subjective judgment.

\subsection{subjective assessment}
\textbf{Applies to:} Based on personal judgment

\textbf{Meaning:} an evaluation partly based on personal interpretation

\textbf{Example:} The diagnosis should not rely entirely on a subjective assessment.

\subsection{subjective interpretation}
\textbf{Applies to:} Based on personal judgment

\textbf{Meaning:} an explanation influenced by one person's perspective

\textbf{Pattern:} a subjective interpretation of + text, event, or evidence

\textbf{Example:} Different historians may offer subjective interpretations of the same event.

\subsection{subjective opinion}
\textbf{Applies to:} Based on personal judgment

\textbf{Meaning:} an opinion based on personal taste or perspective

\textbf{Example:} Whether the new design is attractive is ultimately a subjective opinion.

\textbf{Usage note:} The phrase is common, although opinion is already usually personal.

\subsection{subjective experience}
\textbf{Applies to:} Existing in personal experience

\textbf{Meaning:} an event or condition as it is personally felt or perceived

\textbf{Pattern:} the subjective experience of + noun

\textbf{Example:} The study focused on the subjective experience of chronic pain.

\subsection{subjective perception}
\textbf{Applies to:} Existing in personal experience

\textbf{Meaning:} the way an individual personally understands or senses something

\textbf{Example:} Subjective perception of risk may differ from statistical evidence.

\subsection{subjective well-being}
\textbf{Applies to:} Existing in personal experience

\textbf{Meaning:} a person's own evaluation of happiness and life satisfaction

\textbf{Example:} The survey measured participants' subjective well-being.

\subsection{subjective report}
\textbf{Applies to:} Existing in personal experience

\textbf{Meaning:} a person's description of their own experience

\textbf{Example:} The physician compared the patient's subjective report with the test results.

\subsection{subjective and objective measures}
\textbf{Applies to:} Existing in personal experience

\textbf{Meaning:} self-reported information combined with independently observable data

\textbf{Example:} The researchers used subjective and objective measures of fatigue.

\subsection{reduce subjectivity}
\textbf{Applies to:} Based on personal judgment

\textbf{Meaning:} decrease the influence of personal opinion

\textbf{Example:} Detailed grading criteria can reduce subjectivity in assessment.

\textbf{Usage note:} Subjectivity is the noun form.

\section{Synonyms and antonyms}
\sectionrule
\subsection{Close synonyms}
\subsubsection{personal}
\textbf{Part of speech:} adjective

\textbf{Applies to:} all senses

\textbf{Shared meaning:} Both can refer to an individual's own views or experience.

\textbf{Difference:} Personal is broader and simply means belonging or relating to an individual. Subjective specifically emphasizes that a judgment or experience depends on that individual's perspective.

\textbf{Example:} Her response reflected her personal preferences.

\subsubsection{interpretive}
\textbf{Part of speech:} adjective

\textbf{Applies to:} Based on personal judgment

\textbf{Shared meaning:} Both can describe conclusions shaped by a person's understanding.

\textbf{Difference:} Interpretive describes an explanation based on how evidence or meaning is understood. Subjective more broadly emphasizes personal opinion or feeling.

\textbf{Example:} The exhibition includes an interpretive essay about the artist's work.

\subsubsection{self-reported}
\textbf{Part of speech:} adjective

\textbf{Applies to:} Existing in personal experience

\textbf{Shared meaning:} Both may refer to information based on a person's own account.

\textbf{Difference:} Self-reported specifically means supplied by the person being studied. Subjective is broader and describes the personal nature of the experience or judgment itself.

\textbf{Example:} The study compared self-reported exercise with data from activity monitors.

\subsection{Direct or useful antonyms}
\subsubsection{objective}
\textbf{Part of speech:} adjective

\textbf{Applies to:} all senses

\textbf{Opposition:} Objective judgments and measurements rely on observable facts or agreed standards, while subjective ones depend on personal opinion, perception, or experience.

\textbf{Usage note:} In many real situations, an evaluation can contain both subjective and objective elements.

\textbf{Example:} The researchers used objective measurements of blood pressure.

\section{Words learners may confuse}
\sectionrule
\subsection{biased}
\textbf{Part of speech:} adjective

\textbf{Why learners confuse them:} Both words may describe judgments that are not completely neutral.

\textbf{Key difference:} Subjective means influenced by personal perspective. Biased means unfairly favoring or opposing a person, group, or idea. A judgment can be subjective without being biased.

\textbf{subjective:} The rating is subjective because it depends on personal taste.

\textbf{biased:} The report was biased because it ignored evidence that opposed the author's position.

\subsection{subject to}
\textbf{Part of speech:} adjective versus adjective phrase

\textbf{Why learners confuse them:} The forms begin with the same letters but have unrelated uses.

\textbf{Key difference:} Subjective is an adjective meaning based on personal opinion or experience. Subject to is a separate expression meaning affected by, dependent on, or required to follow something.

\textbf{subjective:} The choice of the best painting is subjective.

\textbf{subject to:} The schedule is subject to change.

\section{Practice exercises}
\sectionrule
\subsection{Word family choice}
Choose the best member of the word family for each blank.

\begin{enumerate}
  \item The detailed scoring guide was introduced to reduce \_\_\_\_\_ in essay grading.
  \item Participants evaluated the taste of each product \_\_\_\_\_.
  \item The survey included both objective data and \_\_\_\_\_ reports.
\end{enumerate}

\subsection{Correct or incorrect?}
Decide whether each sentence is correct. Correct the inaccurate ones.

\begin{enumerate}
  \item Opinions about beauty are often highly subjective.
  \item Because pain is subjective, it is imaginary.
  \item The researchers used objective measurements to reduce subjectivity.
  \item The final results were reported very subjectively.
\end{enumerate}

\subsection{Collocation practice}
Complete each sentence with a natural collocation.

\begin{enumerate}
  \item The choice of the most attractive design is highly \_\_\_\_\_.
  \item The study compared patients' subjective reports \_\_\_\_\_ laboratory measurements.
  \item Clear evaluation criteria can help reduce \_\_\_\_\_ in hiring decisions.
  \item The interviews explored each participant's subjective \_\_\_\_\_ of the event.
\end{enumerate}

\subsection{Complete answer key}
\subsubsection{Word family choice}
\begin{enumerate}
  \item \textbf{subjectivity.} The detailed scoring guide was introduced to reduce subjectivity in essay grading. \emph{Use the noun subjectivity after reduce.}
  \item \textbf{subjectively.} Participants evaluated the taste of each product subjectively. \emph{Use the adverb subjectively to describe how the products were evaluated.}
  \item \textbf{subjective.} The survey included both objective data and subjective reports. \emph{Use the adjective subjective before reports.}
\end{enumerate}

\subsubsection{Correct or incorrect?}
\begin{enumerate}
  \item \textbf{Correct} \emph{Highly subjective correctly describes judgments strongly influenced by personal taste.}
  \item \textbf{Incorrect}. Pain is subjective because it is personally experienced, but it is not imaginary. \emph{Subjective does not mean unreal.}
  \item \textbf{Correct} \emph{Reduce subjectivity is a natural expression for limiting the influence of personal judgment.}
  \item \textbf{Correct} \emph{Subjectively is correctly used as an adverb describing how the results were reported, although the sentence suggests the reporting was influenced by personal perspective.}
\end{enumerate}

\subsubsection{Collocation practice}
\begin{enumerate}
  \item \textbf{subjective.} The choice of the most attractive design is highly subjective. \emph{Highly subjective means strongly dependent on personal taste.}
  \item \textbf{with.} The study compared patients' subjective reports with laboratory measurements. \emph{Subjective reports are often compared with objective measurements.}
  \item \textbf{subjectivity.} Clear evaluation criteria can help reduce subjectivity in hiring decisions. \emph{Reduce subjectivity means decrease the influence of personal opinion.}
  \item \textbf{experience.} The interviews explored each participant's subjective experience of the event. \emph{Subjective experience refers to how an event is personally felt or perceived.}
\end{enumerate}

\vfill
\begin{center}
\small\color{Muted} Created and compiled by \textbf{Amir Kooshky}\\
\href{https://t.me/KooshkyTOEFL}{Telegram Channel}\quad\textbullet\quad
\href{https://www.instagram.com/kooshkytoefl}{Instagram}\quad\textbullet\quad
\href{https://t.me/KooshkyTOEFL\_pv}{Personal Telegram}
\end{center}
\end{document}
